\chapter{Hot reload of UI}
\label{ch:hotreload}

\begin{aside}
Editing UI code and seeing it update without restart is a
huge productivity win. The trick is preserving state across
the reload.
\end{aside}

\section{What hot reload is}

The application binary recompiles. The running process
either:

\begin{itemize}[leftmargin=*]
  \item Swaps the dynamically-loaded UI module.
  \item Restarts but restores state from a snapshot.
\end{itemize}

\section{State-preserving restart}

The simpler approach. On change detection:

\begin{enumerate}[leftmargin=*]
  \item Serialise state to a file.
  \item Exec self with a flag.
  \item New process reads state and resumes.
\end{enumerate}

The user sees a tiny flicker; state is intact.

\section{Module swap}

The complicated approach. The UI lives in a shared library.
On change:

\begin{enumerate}[leftmargin=*]
  \item Compile the new library.
  \item dlopen the new version.
  \item Migrate state (rare: it's all signals; same shape
    usually).
  \item dlclose the old.
\end{enumerate}

This requires careful ABI discipline.

\section{Watching the source}

A file watcher (inotify, FSEvents) detects changes. Debounce
to avoid reloading mid-typing.

\section{Hot reload of just config}

Easier: reload the config file. \maya{} apps that read
config can hot-reload by re-reading and updating signals.

The UI doesn't need to be aware --- signals drive
re-render.

\section{Trade-offs}

\begin{itemize}[leftmargin=*]
  \item State-restart: simpler, slight flicker.
  \item Module swap: smoother but fragile.
\end{itemize}

For development iteration, the simpler is enough.

\section{Recap}

\begin{recap}
\begin{itemize}[leftmargin=*]
  \item Hot reload: state-restart or module swap.
  \item State-restart serialises, execs, resumes.
  \item Module swap dlopens; complex.
  \item Config hot reload is easy and useful.
\end{itemize}
\end{recap}

\section*{Workshop}

\begin{enumerate}[leftmargin=*]
  \item Add config hot reload to your app.
  \item Implement state-restart on source change.
  \item Try module swap; observe the pitfalls.
\end{enumerate}
